\section{构建高性能语音神经假体}

\subsection{实验设置：参与者 T12}

NPTL 实验室在 2022 年招募了一位代号为 \textbf{T12} 的参与者。T12 患有 ALS，虽然还能少量移动手臂，但已经\textbf{无法产生清晰的语音}。研究者在她的大脑中植入了\textbf{四个微电极阵列}：

\begin{itemize}
    \item \textbf{两个阵列}植入\textbf{运动皮层}：预期解码语音产生过程中的口腔面部肌肉运动信号
    \item \textbf{两个阵列}植入 \textbf{Broca 区}：预期解码语言规划信号
\end{itemize}

\begin{knowledgebox}{Broca 区与运动皮层}
\textbf{Broca 区}位于额叶，传统上被认为与语言产生的规划有关。\textbf{运动皮层}负责控制肌肉运动的执行。假设是：Broca 区编码``想说什么''，运动皮层编码``怎么说''。然而实验结果出人意料——后续将详细讨论。
\end{knowledgebox}

\subsection{初步发现：信号主要来自运动皮层}

对 T12 进行行为实验后，研究者得到了一个重要发现：

\begin{itemize}
    \item 运动皮层中的两个阵列能够\textbf{显著高于随机水平}地分类口腔面部运动、单音素和单词
    \item 而 Broca 区中的两个阵列\textbf{几乎无法提供高于随机水平}的分类信息，尤其是在动作执行阶段
\end{itemize}

\begin{warningbox}{Broca 区的``沉默''——出乎意料的结果}
传统神经语言学认为 Broca 区在语言产生中扮演关键角色。然而在本实验中，Broca 区的信号对于解码语音几乎没有帮助。可能的原因包括：（1）电极位置不够精确；（2）Broca 区编码的信息在当前实验范式下不容易被捕获；（3）对 Broca 区功能的传统理解可能需要修正。这一发现仍在进一步研究中。
\end{warningbox}

因此，后续的语音 BCI 系统仅使用运动皮层的两个阵列。

\subsection{数据采集流程}

数据采集是一个需要精心设计的过程：

\begin{enumerate}
    \item 研究者坐在 T12 旁边，屏幕上显示要朗读的句子
    \item T12 尝试发出语音（虽然不清晰），研究者同步记录其神经活动
    \item 形成\textbf{配对数据}：输入是神经活动序列，输出是目标句子的文本
    \item 每个采集 block 包含约 40 个句子，然后休息
    \item 整个采集过程每个研究 session 持续约 100 分钟
    \item 总共采集约 \textbf{10,000 个句子}，来自 Switchboard 电话对话语料库（日常对话英语）
\end{enumerate}

\begin{importantbox}{数据量的限制}
10,000 个句子的训练数据量在机器学习领域算\textbf{非常少}。作为对比，大型语音识别模型通常使用数十万小时的语音数据训练。这一数据限制深刻影响了模型设计的选择——后续会看到为什么不使用 Transformer。
\end{importantbox}

\subsection{系统架构：两阶段解码}

整个语音 BCI 系统采用\textbf{两阶段解码}架构：

\begin{center}
\begin{tikzpicture}[
    box/.style={draw, rounded corners, minimum width=3.5cm, minimum height=1cm, font=\small, align=center},
    arr/.style={-{Stealth[length=3mm]}, thick}
]
\node[box, fill=blue!10] (input) {神经特征序列\\（每 20ms 一帧）};
\node[box, fill=green!10, right=2cm of input] (phone) {音素概率序列};
\node[box, fill=orange!10, right=2cm of phone] (text) {文本输出\\（词序列）};
\draw[arr] (input) -- node[above, font=\scriptsize]{阶段 1} node[below, font=\scriptsize]{GRU + CTC} (phone);
\draw[arr] (phone) -- node[above, font=\scriptsize]{阶段 2} node[below, font=\scriptsize]{Beam Search} node[below=0.4cm, font=\scriptsize]{+ 语言模型} (text);
\end{tikzpicture}
\end{center}

\begin{itemize}
    \item \textbf{阶段 1（神经信号 $\rightarrow$ 音素）}：使用 GRU 神经网络，以 CTC 损失训练，将神经特征序列解码为音素概率序列
    \item \textbf{阶段 2（音素 $\rightarrow$ 文本）}：使用 Beam Search + 语言模型，将音素序列转换为最可能的词序列
\end{itemize}

\subsection{阶段 1：为什么使用 GRU 而不是 Transformer}

\begin{importantbox}{模型选择的理由}
在数据量只有 10,000 句的情况下，模型选择需要考虑以下因素：
\begin{enumerate}
    \item \textbf{数据效率}：Transformer 需要大量数据才能发挥优势；GRU/RNN 在小数据集上表现更好
    \item \textbf{长距离依赖的需求}：语音产生主要涉及\textbf{短程依赖}（相邻音素之间的协调），不需要 Transformer 擅长的超长距离建模
    \item \textbf{实时推理效率}：系统需要在 \textbf{20 毫秒}内完成一次推理；GRU 的计算效率远高于 Transformer，甚至可以在手机上实时运行
\end{enumerate}
\end{importantbox}

最终选择了 \textbf{GRU（Gated Recurrent Unit）}，它是 LSTM 的简化版本，将记忆状态和隐藏状态合并为一个隐藏状态，减少了门控参数，在小数据集上更不容易过拟合。

\subsection{CTC 损失函数}

\begin{knowledgebox}{CTC（Connectionist Temporal Classification）}
CTC 是一种专门为\textbf{输入输出长度不匹配}的序列到序列问题设计的损失函数。它的核心特性包括：

\begin{itemize}
    \item \textbf{长度不匹配处理}：神经特征序列可能有数千帧，而对应的音素序列只有几十个 token
    \item \textbf{单调对齐}：CTC 假设输入输出之间存在单调对齐关系——即前面的输入帧对应前面的输出 token（这与机器翻译的任意对齐不同）
    \item \textbf{空白符}：引入特殊的 blank token 作为``填充''，使输出序列长度与输入一致
    \item \textbf{解码方式}：输出时合并重复 token 并移除 blank，得到最终的 token 序列
\end{itemize}

CTC 最初被广泛应用于\textbf{手写识别}和\textbf{语音识别}（如 CS224S 中介绍的），在语音 BCI 场景中被自然地借鉴过来。
\end{knowledgebox}

CTC 的工作方式可以用一个简单例子说明：假设 GRU 的输出为 \texttt{[-H-HH-EE-LL-LL-OO-]}（其中 \texttt{-} 表示 blank），经过合并重复和移除 blank 后，得到 \texttt{HELLO}。

\subsection{阶段 2：Beam Search 与语言模型}

得到音素概率序列后，需要将其转换为实际的词序列。这一步使用\textbf{修改版的 Beam Search}，结合两种语言模型：

最终的解码目标函数为：

$$\hat{Y} = \arg\max_Y \left[ \log P(Y|X) + \alpha \log P_{\text{LM}}(Y) + \beta |Y| \right]$$

其中：
\begin{itemize}
    \item $P(Y|X)$：GRU + CTC 解码器给出的音素序列概率
    \item $P_{\text{LM}}(Y)$：语言模型给出的句子概率
    \item $|Y|$：句子长度（词数）
    \item $\alpha$：语言模型权重
    \item $\beta$：词插入奖励（word insertion bonus），用于平衡长短句子的概率差异
\end{itemize}

\begin{importantbox}{两层语言模型的设计}
\begin{enumerate}
    \item \textbf{实时解码阶段}使用 \textbf{N-gram 语言模型}：N-gram 模型所有计算都是内存查找，可以在 20ms 内评估上百个候选假设
    \item \textbf{后处理阶段}使用 \textbf{Transformer 语言模型}（如 GPT 级别）：对 Beam Search 产出的 Top-K（如 100 个）候选句子进行\textbf{重排序（reranking）}，选出最终结果。此步骤可以在 $\sim$0.5 秒内完成
\end{enumerate}

这种``先粗后精''的两层策略，既满足了实时性要求，又利用了 Transformer 的强大语言建模能力。
\end{importantbox}

\subsection{系统表现}

最终系统的表现：

\begin{itemize}
    \item 在句子复制任务中，系统可以实时将 T12 的脑信号解码为文字，词错误率约为 \textbf{25\%}
    \item 在问答任务中，系统同样可以解码出 T12 想要表达的回答
    \item \textbf{静默语音}（silent speech）模式——T12 只移动口腔但不发声——系统仍然能够较好地解码
\end{itemize}

\begin{knowledgebox}{静默语音解码的意义}
静默语音解码意味着患者\textbf{不需要实际发出声音}，只需要``尝试''做出说话的口腔动作，系统就能解码。这对于完全丧失发声能力的患者尤为重要——即使声带完全瘫痪，只要运动皮层中仍然编码着语音运动的意图信号，BCI 就有可能解码出来。
\end{knowledgebox}

\subsection{本章小结}

\begin{itemize}
    \item 语音 BCI 采用两阶段架构：GRU + CTC 解码音素，Beam Search + 语言模型生成文本
    \item 在小数据（10,000 句）场景下，GRU 比 Transformer 更合适
    \item CTC 损失函数处理了输入输出长度不匹配和单调对齐的问题
    \item N-gram 语言模型保证实时性，Transformer 语言模型提升最终精度
    \item 系统实现了 $\sim$25\% 词错误率的实时脑到文本解码
\end{itemize}

%% ========== Part 5: 最新进展与未来方向 ==========

